\section{Class 2: External Analysis}
\subsection{Porter Five Forces Framework}
\textbf{Purpose:} Assess industry attractiveness and competitive intensity. Guides
strategic positioning and resource allocation.
\begin{center}
    \begin{tikzpicture}[
            font=\tiny\sffamily,
            >=Latex,
            node distance=3mm and 4mm,
            box/.style={draw, rounded corners=1pt, align=center, minimum width=15mm, minimum height=5mm, fill=white},
            comment/.style={draw, align=left, text width=30mm, rounded corners=1pt, fill=gray!5},
            dashedarrow/.style={-{Latex}, dashed},
            scale=1, transform shape
        ]

        % Central box
        \node[box] (industry) {Industry Rivalry};

        % Surrounding boxes
        \node[box, above=of industry] (new) {Threat of New Entrants};
        \node[box, below=of industry] (sub) {Threat of Substitutes};
        \node[box, left=of industry] (sup) {Bargaining Power of Suppliers};
        \node[box, right=of industry] (buy) {Bargaining Power of Buyers};

        % Arrows
        \draw[->] (new) -- (industry);
        \draw[->] (sub) -- (industry);
        \draw[->] (sup) -- (industry);
        \draw[->] (buy) -- (industry);

    \end{tikzpicture}
\end{center}

\subsubsection{Rivalry between Competitors}
\begin{itemize}
    \item \textbf{Price Competition:} Intense when products are undifferentiated and switching
          costs are low. Cost alone determines which product customers buy and the most
          efficient producer wins.
    \item \textbf{Differentiation:} Intense when products are similar and switching costs are
          low. Firms compete on non-price factors (e.g., quality, brand). Customer
          loyalty can reduce price sensitivity and allow for higher margins.
    \item \textbf{Focus on Niche:} Intense when firms target specific market segments. Firms can
          compete on price or differentiation within the niche. Success depends on how
          wellthe firm understands and serves the niche's unique needs.
\end{itemize}

\textbf{Generic Strategies:} Cost Leadership, Differentiation, Focus. Firms must choose
one to avoid being ``stuck in the middle'' and losing competitive advantage.

\subsubsection{Bargaining power of Suppliers}

\textbf{High Power:} Few suppliers, unique products, high switching costs. Can demand
higher prices or reduce quality. Firms may need to integrate backward or find
alternative suppliers.

\textbf{Low Power:} Many suppliers, standardized products, low switching costs. Firms
can negotiate better terms and have more flexibility. Focus on cost control and
supplier relationships.

\textbf{Examples:} High power Airbus and Boeing for commercial aviation; Low power -
Grab Drivers providing service for Grab.

\subsubsection{Bargaining power of Buyers}
\textbf{High Power:} Few buyers, standardized products, low switching costs, low search
cost, low transportation cost. Can demand lower prices, better quality, or
additional services. Firms may need to differentiate or focus on customer
loyalty.

\textbf{Low Power:} Many buyers, unique products, high switching costs. Firms have more
pricing power and can focus on product development and marketing.

\textbf{Examples:} High power - Travel Agencies can pick from many airlines to fly to
holiday destinations; Low power - Individual consumers buying smartphones.

\subsubsection{Threat of New Entrants}
\textbf{High Threat:} Low barriers to entry, high industry growth, low capital
requirements, easy access to distribution channels, little goverement
regulation. New entrants can quickly gain market share and intensify
competition. Incumbents may need to innovate or lower prices to maintain their
position.

\textbf{Low Threat:} High barriers to entry, slow industry growth, high capital
requirements, difficult access to distribution channels, significant government
regulation. Incumbents have a strong position and can focus on long-term
strategies without worrying about new competitors.

\textbf{Examples:} High threat - Food and Beverage industry with low capital
requirements; Low threat - Pharmaceuticals, where many phases of trials and
government regulations have to be passed.

\subsubsection{Threat of Substitutes}
\textbf{High Threat:} Many substitutes, low switching costs, attractive
price-performance ratio. Can limit industry profitability as customers switch
to alternatives. Firms may need to differentiate or innovate to stay
competitive.

\textbf{Low Threat:} Few substitutes, high switching costs, poor price-performance
ratio. Firms have more pricing power and can focus on improving their products
and services.

\textbf{Examples:} High threat Public transportation as a substitute for private car
ownership; Low threat - Electricity as a substitute for gasoline in the short
term.

\subsection{SWOT Analysis}
\textbf{Purpose:} Identify internal strengths and weaknesses, and external
opportunities and threats. Guides strategic planning and decision-making.
\begin{tabularx}{\linewidth}{|X|X|}
    \hline
    \rowcolor{gray!20} \textbf{Strengths} (Internal, +)     & \textbf{Weaknesses} (Internal, -)                             \\ \hline
    \textit{What do we do well?}                            & \textit{Where do we lack resources?}                          \\ e.g., Brand equity (Starbucks), Hub dominance
    (Airlines).                                             & e.g., High fixed costs, rigid labor unions.                   \\ \hline
    \rowcolor{gray!20} \textbf{Opportunities} (External, +) & \textbf{Threats} (External, -)                                \\ \hline
    \textit{External factors to exploit.}                   & \textit{External factors that risk profits.}                  \\ e.g., Emerging markets (China), falling
    oil prices.                                             & e.g., New LCC entrants, health trends, rising interest rates. \\
    \hline
\end{tabularx}

\textbf{Strategy (TOWS Matrix):}
\begin{itemize}
    \item \textbf{S-T:} Use strengths to mitigate threats.
    \item \textbf{W-O:} Address weaknesses to capitalize on opportunities.
    \item \textbf{S-O:} Use strengths to exploit opportunities.
    \item \textbf{W-T:} Defensive strategy to minimize weaknesses and avoid threats.
\end{itemize}

\section{Case: IBM's Strategic Pivot (1993)}
\textbf{Context:} IBM faced an existential crisis as the industry shifted from a \textbf{Vertical} (proprietary/integrated) model to a \textbf{Horizontal} (fragmented/specialized) model.

\subsection{Industry Shift: Vertical vs. Horizontal}
\begin{itemize}
    \item \textbf{Vertical (Old):} One firm (IBM) provided the processor, OS, and apps. High
          margins but slow innovation.
    \item \textbf{Horizontal (New):} Niche players (Intel, Microsoft, Oracle) dominated specific
          ``slices'' of the stack.
    \item \textbf{The Integration Crisis:} Customers struggled to integrate fragmented
          ``piece-parts'' from different vendors, creating a ``cacophony of standards.''
\end{itemize}

\subsection{IBM SWOT Analysis (1993)}
\begin{tabularx}{\linewidth}{|X|X|}
    \hline
    \rowcolor{gray!20} \textbf{Strengths}        & \textbf{Weaknesses} \\ \hline
    \begin{itemize}
        \item Huge Scale/Broad Capabilities
        \item World-Class R\&D (Nobel labs)
        \item Deeply entrenched customer roots
    \end{itemize}   &
    \begin{itemize}
        \item Inefficient Cost Structure (42\textcent/\$1 rev)
        \item Bureaucratic ``Check'' Culture
        \item Mainframe Revenue Dependence
    \end{itemize}          \\ \hline
    \rowcolor{gray!20} \textbf{Opportunities}    & \textbf{Threats}    \\ \hline
    \begin{itemize}
        \item Demand for ``Systems Integrators''
        \item IT Services Growth
        \item Rise of Networked Computing
    \end{itemize} &
    \begin{itemize}
        \item ``Wintel'' (MSFT/Intel) Duopoly
        \item Commodity Pricing/PC Revolution
        \item Commercialization Failures (UNIX)
    \end{itemize}                         \\ \hline
\end{tabularx}

\subsection{The Gerstner Strategy (The Turnaround)}
\begin{itemize}
    \item \textbf{Reject the Breakup:} Leveraged IBM's size as a strength. IBM became the \textbf{Integrator} to solve customer complexity.
    \item \textbf{Services as the ``Tail'':} Shifted focus from selling hardware to providing
          end-to-end IT services and consulting.
    \item \textbf{Economic Right-Sizing:} Cut \$8.9B in expenses and 35k jobs to close the
          42\textcent\ vs 31\textcent\ efficiency gap.
    \item \textbf{Stabilize the Core:} Lowered mainframe prices to prevent customer defection,
          sacrificing short-term margin for long-term survival.
\end{itemize}

\textbf{Strategic Takeaway:} IBM survived by pivoting from a product-centric firm to a \textbf{solutions-centric} firm, monetizing the very ``complexity'' that the fragmented
industry created.